% BHU PYQ -- Manually transcribed & error-corrected
% Paper: BCH-213 : Specialised Accounts - I
% Exam: B.Com. (Hons.) Semester III Examination, 2013-14

\documentclass[11pt,a4paper]{article}
\usepackage[a4paper,top=2.5cm,bottom=2.5cm,left=2.8cm,right=2.8cm,headheight=14pt]{geometry}
\usepackage{amsmath,amssymb}
\usepackage[T1]{fontenc}
\usepackage[utf8]{inputenc}
\usepackage{lmodern,microtype,fancyhdr,enumitem,hyperref,lastpage,parskip}

\hypersetup{
    colorlinks=true,
    urlcolor=black,
    linkcolor=black,
    pdftitle={BCH-213: Specialised Accounts - I -- BHU B.Com. (Hons.) Sem III 2013-14}
}

\pagestyle{fancy}
\fancyhf{}
\renewcommand{\headrulewidth}{0.4pt}
\renewcommand{\footrulewidth}{0.4pt}
\fancyhead[L]{\small\textit{Banaras Hindu University}}
\fancyhead[R]{\small\textit{BCH-213: Specialised Accounts - I}}
\fancyfoot[C]{\small Page \thepage\ of \pageref{LastPage}}

\newlist{parts}{enumerate}{2}
\setlist[parts,1]{label=(\alph*),leftmargin=*,itemsep=5pt,topsep=3pt}
\setlist[parts,2]{label=(\roman*),leftmargin=*,itemsep=3pt,topsep=2pt}
\newcommand{\pts}[1]{\hfill\textit{[#1]}}

\begin{document}

\begin{center}
  {\large\bfseries BANARAS HINDU UNIVERSITY}\\[2pt]
  {\normalsize B.Com. (Hons.) --- Semester III Examination, 2013-14}\\[6pt]
  \rule{0.8\linewidth}{0.5pt}\\[6pt]
  {\large\bfseries Paper No. BCH-213: Specialised Accounts - I}\\[4pt]
  {\normalsize Time: 3 Hours\hfill Maximum Marks: 70}
\end{center}

\rule{\linewidth}{0.4pt}

\smallskip
\noindent\textit{Instructions: Answer all questions. All questions carry equal marks (14 marks each).}

\bigskip

\noindent\textbf{Question 1.} \\
What do you understand by Royalty? Explain the various types of Royalty.\pts{14}

\centerline{\textbf{OR}}

A Coal Company takes a lease of coal mine for a term of 8 years from 1st January, 2010 on a minimum rent of Rs. 12,000 per year merging into a royalty of 75 paise per ton of coal raised payable half-yearly on 30th June and 31st December each year. Shortworking can be recouped during the first three years of lease only. Royalty which was due on 31st December, 2011 was paid on 2nd January, 2012. Coal Company closes its books of accounts on 31st December each year. Coal raised is as follows:\pts{14}
\begin{center}
\renewcommand{\arraystretch}{1.2}
\begin{tabular}{|l|r|}
\hline
\textbf{Half-year ending} & \textbf{Coal raised (in tons)} \\ \hline
30th June, 2010 & 6,000 \\
31st December, 2010 & 8,000 \\
30th June, 2011 & 10,000 \\
31st December, 2011 & 6,000 \\
30th June, 2012 & 7,000 \\
31st December, 2012 & 9,000 \\ \hline
\end{tabular}
\end{center}
Prepare necessary Accounts in the books of Coal Company.

\medskip\rule{\linewidth}{0.2pt}

\noindent\textbf{Question 2.} \\
\begin{parts}
  \item Distinguish between Hire-purchase System and Instalment System.\pts{7}
  \item Show imaginary Journal Entries in the books of a purchaser in case of Hire-purchase System.\pts{7}
\end{parts}

\centerline{\textbf{OR}}

Shyam purchased a machine on 1st January, 2010 on the Hire-purchase System from Sohan. The cash price of the machine was Rs. 2,23,500 and the payment was to be made as follows: Rs. 60,000 on signing of agreement and the balance in three instalments of Rs. 60,000 each at the end of the year. 5\% interest is charged by the seller per annum. Shyam has decided to write off 10\% depreciation annually on the diminishing balance of the value. The purchaser could not pay off the instalment due on 31st December, 2011, and as the hire-purchase price had not been paid, the seller after completing all formalities, took possession of the machine.\pts{14}

Prepare necessary Accounts in the books of Shyam and Sohan.

\medskip\rule{\linewidth}{0.2pt}

\noindent\textbf{Question 3.} \\
Point out different types of branches for accounting purposes. Describe the rules relating to conversion of Foreign Branch Trial Balance.\pts{14}

\centerline{\textbf{OR}}

M Ltd. has head office at Mumbai and branches in other cities. It supplies goods to branches at cost plus 25\%. Accounts are kept at the head office from where all expenses, except petty expenses are paid. Branches are allowed to maintain petty cash balance of Rs. 400 on imprest system. From the following balances, prepare Branch A/c assuming that accounts are closed on 31st December each year:\pts{14}

\textbf{Branch Balances on 1.1.2012:}
\begin{center}
\renewcommand{\arraystretch}{1.2}
\begin{tabular}{|l|r|}
\hline
\textbf{Particulars} & \textbf{Amount (Rs.)} \\ \hline
Petty cash & 400 \\
Stock at selling price & 20,000 \\
Sundry debtors & 4,000 \\
Sundry creditors & 1,200 \\
Furniture & 10,000 \\
Rent paid up to 31.03.2013 & 300 \\
Wages outstanding & 100 \\ \hline
\end{tabular}
\end{center}

\textbf{Transactions for the year 2012:}
\begin{center}
\renewcommand{\arraystretch}{1.2}
\begin{tabular}{|l|r|}
\hline
\textbf{Particulars} & \textbf{Amount (Rs.)} \\ \hline
Goods sent to branch less returns & 1,00,000 \\
Cash sales & 80,000 \\
Credit sales & 45,000 \\
Allowances to debtors & 500 \\
Cash received from customers & 40,000 \\
Bad debts & 20 \\
Discount & 80 \\
Cash purchases by branch on permission from H.O. & 10,600 \\
Cash paid to creditors & 8,000 \\ \hline
\end{tabular}
\end{center}

\textbf{Payments made by H.O.:}
\begin{center}
\renewcommand{\arraystretch}{1.2}
\begin{tabular}{|l|r|}
\hline
\textbf{Particulars} & \textbf{Amount (Rs.)} \\ \hline
Rent for one year paid on 1.4.2012 & 1,800 \\
Salaries up to November, 2012 & 2,200 \\
Insurance paid for one year on 1.4.2012 & 360 \\
Wages & 100 \\ \hline
\end{tabular}
\end{center}

\textbf{Payments made by Branch:}
\begin{center}
\renewcommand{\arraystretch}{1.2}
\begin{tabular}{|l|r|}
\hline
\textbf{Particulars} & \textbf{Amount (Rs.)} \\ \hline
Petty expenses & 180 \\
Other expenses & 500 \\ \hline
\end{tabular}
\end{center}

\textbf{Balances on 31.12.2012:}
\begin{center}
\renewcommand{\arraystretch}{1.2}
\begin{tabular}{|l|r|}
\hline
\textbf{Particulars} & \textbf{Amount (Rs.)} \\ \hline
Creditors & 3,000 \\
Stock at cost & 30,000 \\ \hline
\end{tabular}
\end{center}
Charge 10\% depreciation on furniture.

\medskip\rule{\linewidth}{0.2pt}

\noindent\textbf{Question 4.} \\
What are the objectives of Departmental Accounts? What are the bases of apportioning expenses among different departments?\pts{14}

\centerline{\textbf{OR}}

The directors of a Departmental Stores Ltd. wish to ascertain approximately the net profits of the A, B and C departments separately for the quarter ended March 31st, 2013. It is found impracticable actually to take stock on that date but an adequate system of department accounts is in use and the normal rates of gross profit for the departments concerned are 40\%, 30\% and 20\% on turnover respectively. Indirect expenses are charged in proportion to departmental turnover. Following are the figures for each department:\pts{14}
\begin{center}
\renewcommand{\arraystretch}{1.2}
\begin{tabular}{|l|r|r|r|}
\hline
\textbf{Particulars} & \textbf{Dept. A (Rs.)} & \textbf{Dept. B (Rs.)} & \textbf{Dept. C (Rs.)} \\ \hline
Stock on 1.1.2013 & 30,000 & 35,000 & 15,000 \\
Purchases to March 31, 2013 & 35,000 & 37,500 & 23,500 \\
Sales to March 31, 2013 & 60,000 & 50,000 & 30,000 \\
Direct expenses & 10,100 & 7,250 & 3,550 \\ \hline
\end{tabular}
\end{center}
Total indirect expenses for the period (including those relating to other departments) were Rs. 21,000 on total sales of Rs. 4,20,000.

Prepare a Statement showing gross profit, net profit after making reserve for stock at 10\% in respect of each department.

\medskip\rule{\linewidth}{0.2pt}

\noindent\textbf{Question 5.} \\
Discuss the meaning and objectives of HR Accounting.\pts{14}

\centerline{\textbf{OR}}

Discuss the present value approaches to HR Accounting.\pts{14}

\vfill
\rule{\linewidth}{0.4pt}
\begin{center}\small
\href{https://bhu-exam.vercel.app/}{Click here to visit our website}\\[4pt]
\href{https://me-aryan.vercel.app/}{Click here to visit our contributor}\\[4pt]
\href{https://quashan-me.vercel.app/}{Click here to visit our contributor}
\end{center}

\end{document}
